% Referencias · Se incluye desde informe.tex.
\begin{thebibliography}{}
% Formato APA 7, adaptación en español: «y» entre autores, «s. f.», «Recuperado el ... de».
% Orden alfabético por primer autor. Etiqueta natbib: [cita corta(año)].
% Verificación de cada referencia: docs/verificacion_referencias.md (issue #24).

\bibitem[Campo~Agredo(2017)]{campo2017}
Campo Agredo, C. A. (2017). \emph{Desarrollo de modelos de predicción sobre el hurto a personas en Bogotá a partir del análisis y aplicación de algoritmos en base de datos} [Trabajo de grado de pregrado, Universidad de los Andes]. Repositorio Institucional Séneca. \url{https://hdl.handle.net/1992/50785}

\bibitem[{Fiscalía General de la Nación}(2026)]{fiscalia}
Fiscalía General de la Nación. (2026). \emph{Procesos Fiscalía V3} [Conjunto de datos; corte del 31 de octubre de 2026]. Portal de Datos Abiertos de Colombia. Recuperado el 11 de noviembre de 2026, de \url{https://www.datos.gov.co/d/dbdv-iihs}

\bibitem[Fonseca~Vásquez(2019)]{fonseca2019}
Fonseca Vásquez, A. R. (2019). \emph{Análisis espacial del hurto a personas para la toma de decisiones del Comando de Atención Inmediata (CAI) en la localidad de Los Mártires Bogotá D. C. entre los años 2017--2018} [Trabajo de especialización, Universidad Militar Nueva Granada]. Repositorio Institucional UMNG. \url{https://repository.umng.edu.co/server/api/core/bitstreams/2e3a1b76-45f4-4348-82e1-ace03d817513/content}

\bibitem[Garzón~Acosta(2025)]{garzon2025}
Garzón Acosta, J. J. (2025). \emph{Modelos de predicción del hurto a personas en Bogotá D. C. (2015--2024): Un enfoque comparativo entre MCE y VECM} [Trabajo de especialización, Fundación Universitaria Los Libertadores]. Repositorio Institucional Los Libertadores. \url{https://repository.libertadores.edu.co/bitstreams/45b1a97c-3c66-4beb-988a-95e7dfb9c14b/download}

\bibitem[Martínez~Chitiva(2025)]{martinez2025}
Martínez Chitiva, A. V. (2025). \emph{Análisis espacio-temporal para caracterización, predicción y mitigación del hurto a personas en Bogotá} [Tesis de maestría, Universidad de Manizales]. Repositorio Institucional RIDUM. \url{https://ridum.umanizales.edu.co/handle/20.500.12746/7648}

\bibitem[{Policía Nacional de Colombia}(2026)]{policia}
% Recurso usado (D-01, issue #18): archivos anuales 2020--2025 a nivel de registro.
Policía Nacional de Colombia. (2026). \emph{Información de delitos a nivel de registro} [Conjunto de datos; archivos anuales 2020--2025]. Dirección de Investigación Criminal e INTERPOL. Recuperado el 11 de noviembre de 2026, de \url{https://www.policia.gov.co/estadistica-delictiva}

\bibitem[Provost y Fawcett(2013)]{provost2013}
Provost, F., y Fawcett, T. (2013). \emph{Data science for business: What you need to know about data mining and data-analytic thinking}. O'Reilly Media.

\bibitem[Riascos et~al.(2020)]{riascos2020}
% Número confirmado (issue #24): el listado oficial de la Escuela de Gobierno la presenta como
% «Nota de política No. 38». El «34» de los metadatos del PDF es el nombre del archivo de diseño.
Riascos, A. J., Dulce, M., Moreno, J. S., y Gómez, F. (2020). \emph{Prediciendo el crimen en Bogotá} (Nota de Política n.\textsuperscript{o}~38). Escuela de Gobierno Alberto Lleras Camargo, Universidad de los Andes. \url{https://gobierno.uniandes.edu.co/wp-content/uploads/NP-38.pdf}

\bibitem[Riascos~Villegas et~al.(2023)]{riascos2023}
Riascos Villegas, Á. J., Ñungo, J. S., Gómez Tobón, L., Dulce Rubio, M., y Gómez, F. (2023). Modelling underreported spatio-temporal crime events. \emph{PLOS ONE}, \emph{18}(7), Artículo e0287776. \url{https://doi.org/10.1371/journal.pone.0287776}

\bibitem[{Secretaría Distrital de Seguridad, Convivencia y Justicia}(2026)]{linea123}
Secretaría Distrital de Seguridad, Convivencia y Justicia. (2026). \emph{Llamadas tramitadas C4 Bogotá NUSE Línea 123, enero 2015 -- septiembre 2026} [Conjunto de datos]. Datos Abiertos Bogotá. Recuperado el 11 de noviembre de 2026, de \url{https://datosabiertos.bogota.gov.co/dataset/incidentes}

\bibitem[Urcuqui et~al.(2020)]{urcuqui2020}
Urcuqui, C., Moreno, J., Montenegro, C., Riascos, A., y Dulce, M. (2020). Accuracy and fairness in a conditional generative adversarial model of crime prediction. En \emph{2020 7th International Conference on Behavioural and Social Computing (BESC)} (pp. 1--6). IEEE. \url{https://doi.org/10.1109/BESC51023.2020.9348315}

\end{thebibliography}
